\section{Introduction}\label{sec:introduction}

LLT polynomials were introduced by A.~Lascoux, B.~Leclerc and J.-Y.~Thibon as
$q$-deformations of products of Schur functions~\cite{LLT1997}.
I.~Grojnowski and M.~Haiman proved that they are Schur-positive
\cite{GrojnowskiHaiman2006}.
LLT polynomials are the building blocks of the combinatorial formula for
Macdonald polynomials of J.~Haglund, M.~Haiman and
N.~Loehr~\cite{HaglundHaimanLoehr2005}, and of the Shuffle Theorem, conjectured
in~\cite{HaglundHaimanLoehrRemmelUlyanov2005} and proved by E.~Carlsson and
A.~Mellit~\cite{CarlssonMellit2017}; see also~\cite{qtCatalanBook}.

The vertical-strip LLT polynomials form a subfamily which can be described by
colorings of decorated unit-interval graphs~\cite{AlexanderssonPanova2018}: if
$\Gamma$ is such a graph on $[n]$, then
\[
 \LLT_\Gamma(\xvec;q)=\sum_\kappa q^{\asc(\kappa)}x_{\kappa(1)}x_{\kappa(2)}\dotsm x_{\kappa(n)},
\]
where $\kappa$ ranges over the colorings which increase strictly along the
strict edges, and $\asc(\kappa)$ counts the other edges along which $\kappa$
increases.
Without strict edges, they are the unicellular LLT polynomials, which are
related by a plethystic substitution to the chromatic quasisymmetric functions of
J.~Shareshian and M.~Wachs~\cite{ShareshianWachs2012,ShareshianWachs2016},
see~\cite{CarlssonMellit2017,GuayPaquet2016x}.
For chromatic symmetric functions, elementary positivity is the content of the
Stanley--Stembridge conjecture~\cite{StanleyStembridge1993,Stanley1995}.
M.~Guay-Paquet reduced it to unit-interval graphs~\cite{GuayPaquet2013}, and
T.~Hikita proved it~\cite{Hikita2024}.
The refinement of Shareshian and Wachs to general $q$~\cite{ShareshianWachs2016}
remains open.

The first named author and R.~Sulzgruber proved that, after the substitution
$q\mapsto q+1$, the vertical-strip LLT polynomials expand positively in the
elementary symmetric functions~\cite{AlexanderssonSulzgruber2022}, as conjectured
in~\cite{GarsiaHaglundQiuRomero2019,Alexandersson2019llt}; the positivity was
obtained independently by M.~D'Adderio~\cite{DAdderio2020}.
The expansion is
\begin{equation}\label{eq:introAS}
 \LLT_\Gamma(\xvec;q+1)=\sum_{\theta}q^{\asc(\theta)}\elementaryE_{\lambda(\theta)}(\xvec).
\end{equation}
Here $\theta$ ranges over the orientations of the normal edges of $\Gamma$,
$\asc(\theta)$ is the number of ascending edges, those oriented from the smaller
to the larger vertex, and $\lambda(\theta)$ records the sizes of the sets of vertices with the same highest reachable vertex.
The proof in~\cite{AlexanderssonSulzgruber2022} shows that both sides
of~\eqref{eq:introAS} satisfy a system of linear relations which determines them.
Local linear relations of this kind go back to the modular law of
M.~Guay-Paquet~\cite{GuayPaquet2013} and to the relations of
S.~J.~Lee~\cite{Lee2018x}; see also~\cite{AbreuNigro2020ax}.
The analogous question for circular unit-interval graphs, whose chromatic
quasisymmetric functions were studied by B.~Ellzey~\cite{Ellzey2016,Ellzey2017},
was left open in~\cite[Sec.~6.8]{AlexanderssonSulzgruber2022}.

The same polynomial has a simpler expansion.
Expanding $(1+q)^{\asc(\kappa)}$ shows that
\begin{equation}\label{eq:introK}
 \LLT_\Gamma(\xvec;q+1)=\sum_{\theta}q^{\asc(\theta)}\Kp_{\theta}(\xvec),
\end{equation}
where $\Kp_\theta$ is the generating function of strict $P$-partitions for the
partial order $P=\Pplus{\theta}$ generated by the strict edges and the ascending
edges of $\theta$.
The summands $\Kp_{\theta}$ are not symmetric functions.
The situation is the same for chromatic symmetric functions, which are sums of
strict $P$-partition generating functions over the acyclic orientations of a
graph~\cite{Stanley1995}.
B.-H.~Hwang introduced flips of acyclic orientations, which reverse a fork of two
covers, and showed that the corresponding sums over flip classes are
symmetric~\cite[Thm.~3.8]{Hwang2022x} and Schur-positive.
He conjectured that they are elementary-positive~\cite[Conj.~3.10]{Hwang2022x}.
He also mentioned an analogous flip on orientations for LLT polynomials,
suggested by the first named author, and asked whether its class sums are
Schur-positive~\cite[Sec.~5.3]{Hwang2022x}.
This is the flip of the present paper.

\subsection*{Main results}

Let $u<v<w$ and let $\theta$ be an orientation in which $uv$ is ascending and
$vw$ is not.
If $uv$ and $uw$ are covers of $\Pplus{\theta}$, and $vw$ and $uw$ are covers
after reversing $uv$ and $vw$, we say that the two orientations differ by a
\emph{flip}, see Definition~\ref{def:flip}.
Thus a flip keeps a long edge $uw$ and moves an ascending short edge across the
triangle.
A flip preserves $\asc(\theta)$, and it changes $\Kp_{\theta}$ as well as
$\lambda(\theta)$.

\begin{theoremA}[Theorem~\ref{thm:main}]
Let $\Gamma$ be an ordinary or circular decorated unit-interval graph, and let
$C$ be a flip class of orientations of $\Gamma$.
Then
\[
 \sum_{\theta\in C}\Kp_{\theta}(\xvec)=\sum_{\theta\in C}\elementaryE_{\lambda(\theta)}(\xvec).
\]
\end{theoremA}

For example, in the complete graph on three vertices the two orientations with
ascending edges $\{12,13\}$ and $\{13,23\}$ form a flip class, and
\[
 \Kp_{\{12,13\}}+\Kp_{\{13,23\}}=\elementaryE_{21}+\elementaryE_3,
\]
see Example~\ref{ex:smallestFlip}.
Thus every class sum is elementary-positive, and in particular Schur-positive,
which answers the question of Hwang mentioned above.
Summing Theorem~A over all flip classes and comparing~\eqref{eq:introK}
with~\eqref{eq:introAS} we obtain the following.

\begin{theoremB}[Corollary~\ref{cor:circularExpansion}]
The expansion~\eqref{eq:introAS} holds for every circular decorated unit-interval
graph.
Here $\theta$ ranges over the admissible orientations, in which the strict and
ascending edges form no directed cycle, and highest reachable vertices are
computed on the lift to $\setZ$.
\end{theoremB}

For ordinary graphs this is~\eqref{eq:introAS} itself, and our proof does not use
the recursion of~\cite{AlexanderssonSulzgruber2022}.
Theorem~A refines~\eqref{eq:introAS}: every flip class contributes an
elementary-positive summand $q^{\asc(C)}\sum_{\theta\in C}\elementaryE_{\lambda(\theta)}$.
The same refinement applies to the LLT cumulants of M.~Dołęga and
M.~Kowalski~\cite{DolegaKowalski2021x}.

\begin{corollaryC}[Corollary~\ref{cor:cumulant}]
The LLT cumulants of circular decorated unit-interval graphs, after
$q\mapsto q+1$, are sums of elementary-positive contributions of flip classes.
\end{corollaryC}

\subsection*{Method}

The proof of Theorem~A for ordinary graphs uses one local identity: for three
vertices $u<v<w$ and any fixed set of other edges containing $uw$, the four ways
of adding the short edges $uv$ and $vw$ satisfy
\[
 \Psi(\text{with } uv)+\Psi(\text{with } vw)=\Psi(\text{with neither})+\Psi(\text{with both}),
\]
both for $\Psi=\Kp$ and for $\Psi=\elementaryE_\lambda$, see Lemma~\ref{lem:fourTerm}.
Summed over all orientations with a fixed number of edges, it gives refined
forms of the two basic bounce relations of~\cite{AlexanderssonSulzgruber2022}.
To pass from all orientations to a single flip class, we prove a connectivity
theorem: in a graph without articulations, all connected orientations with a
given number of ascending edges in which every corner is a cover form one flip
class.
Every other flip class of that size is expressed, with the same coefficients for
both sides, in terms of flip classes of graphs with fewer edges.
A graph with articulations is cut at the end of its first block, and the
relation is transported to the whole graph by fixing the color of the cut vertex.
The circular case is reduced to the ordinary one: a finite signed expansion
writes every admissible orientation in terms of orientations whose edge lengths are
differences of a potential, the coefficients are constant on flip classes, and
the flip classes of such orientations can be drawn on a line.

Theorem~A for ordinary and circular graphs, and the elementary expansions
obtained from it, have been formally verified in the Lean proof assistant using
Mathlib; the formalization is available at
\url{https://github.com/PerAlexandersson/VerticalStripLLT}.

\subsection*{Outline}

Section~\ref{sec:preliminaries} contains the definitions and the precise
statements.
Section~\ref{sec:local} proves the four-term identity.
Section~\ref{sec:ordinary} proves Theorem~A for ordinary graphs, assuming the
connectivity theorem, which is proved in Section~\ref{sec:connectivity}.
Section~\ref{sec:circular} treats circular graphs, and
Section~\ref{sec:consequences} derives the elementary expansion and the result
on cumulants, and ends with two questions.
